% 共享排版样式（对标《工程本体论》handbook：三盒子 + 中文 + 等宽）
% 《产品可信工程》候选样书元数据的唯一排版来源。
% 正式出版前，所有“待……”字段必须由作者与出版单位确认并替换；不得据此生成 ISBN/CIP。
\newcommand{\BookEnglishTitle}{PRODUCT TRUSTWORTHINESS}
\newcommand{\BookTitle}{产品可信工程}
\newcommand{\BookSubtitle}{人读规范与 Semantica 唯一执行语义}
\newcommand{\BookScopeLine}{以 ISO 26262 为贯穿样板，以 Semantica 为唯一机器执行入口}
\newcommand{\BookVolumeLabel}{第二卷}
\newcommand{\BookVolumeName}{规范语义与 Semantica 深度绑定}
\newcommand{\BookEditionLabel}{Semantica 深度绑定候选审阅版}
\newcommand{\BookVersion}{2026.08-semantica.1}
\newcommand{\BookDateCN}{2026 年 8 月}
\newcommand{\BookSemanticContractLine}{书是供人阅读、评审和追责的规范；Semantica 是全书唯一可执行语义。}
\newcommand{\BookSemanticaStatus}{章节包的可执行状态以 Semantica 注册表、验收收据与发布记录为准；本候选书不宣称全部章节包已经发布完成。}

% 内部文档号仅供候选稿追踪，绝不是 ISBN、CIP 核字号或出版备案号。
\newcommand{\BookInternalID}{PTW-V2-2026.08-SEMANTICA.1}
\newcommand{\BookCoverCredit}{claude F5 max}
\newcommand{\BookCoverWorkNote}{5k · 用时 4 天}
\newcommand{\BookAuthorStatus}{claude F5 max}
\newcommand{\BookRightsHolderStatus}{待版权权利人确认}
\newcommand{\BookProductionRolesStatus}{责任编辑、责任校对与装帧设计待出版流程确认}
\newcommand{\BookSeriesStatus}{两册候选书之第二卷；正式丛书名待出版方案确认}
\newcommand{\BookISBNStatus}{待出版社分配（本候选版无 ISBN）}
\newcommand{\BookCIPStatus}{待出版申报（本候选版无 CIP 数据）}
\newcommand{\BookPublisherStatus}{待出版决定}
\newcommand{\BookEditionStatus}{Semantica 深度绑定候选审阅版；正式版次、印次及章节包发布状态分别决定}
\newcommand{\BookPrintStatus}{候选电子样书为 A4；正式开本、出血、装订、印张与字数待定}
\newcommand{\BookPriceStatus}{待出版决定}
\newcommand{\BookProjectURL}{https://github.com/jiaqiwang969/OntologyEngineering}

\usepackage[a4paper,margin=2.6cm]{geometry}
\usepackage{booktabs,tabularx,xltabular,amssymb,amsmath}
\usepackage{graphicx}
\usepackage{adjustbox}   % 图宽自适应：max size 封顶不放大
\usepackage{rotating}    % 极宽图 sidewaysfigure 横排
\usepackage[font=small,labelfont=bf,labelsep=quad,hypcap=false]{caption}
\usepackage{xcolor}
\usepackage[most]{tcolorbox}
\usepackage{tikz}
\usetikzlibrary{calc,fadings}
\usepackage{fontspec}
\usepackage{needspace}
\usepackage{emptypage}
\usepackage{fvextra}
\usepackage[topmarks]{titleps}
\setmonofont{Menlo}[Scale=0.86]
\usepackage{xurl}
\usepackage{hyperref}
\urlstyle{tt}
\definecolor{bookslate}{RGB}{54,78,99}
\definecolor{bookteal}{RGB}{45,139,139}
\definecolor{bookamber}{RGB}{214,145,52}
\definecolor{booknight}{RGB}{7,20,33}
\definecolor{bookice}{RGB}{178,220,240}
\definecolor{bookcoral}{RGB}{239,82,54}
\hypersetup{
  hidelinks,
  pdftitle={\BookTitle：\BookSubtitle},
  pdfauthor={\BookAuthorStatus},
  pdfsubject={\BookScopeLine；书是人读规范，Semantica 是唯一执行语义},
  pdfkeywords={产品可信工程, ISO 26262, 功能安全, 工程本体, Semantica, 唯一执行语义, 证据闭环},
  bookmarksnumbered=true,
  unicode=true
}

% 中文正文避免段首/段尾单行跨页；代码框前导句由构建器用 \nopagebreak 绑定。
\clubpenalty=10000
\widowpenalty=10000
\displaywidowpenalty=10000
\brokenpenalty=10000
\setlength{\emergencystretch}{3.5em}
\raggedbottom

% 成书页眉页脚：运行页眉镜像排布；页码置于外侧页脚。
% 内部文档号只保留在内书名页，不进正文页脚（读者视图保持安静）。
% 章、部和目录首页使用安静的 plain 样式；openright 自动空白页仍由 emptypage 清空。
\newpagestyle{bookfront}[\small\sffamily\color{bookslate!82!black}]{
  \sethead[\BookTitle][][\BookVolumeLabel\quad\BookVolumeName]
          {\BookVolumeLabel\quad\BookVolumeName}{}{\BookTitle}
  \setfoot[\thepage][][]
          {}{}{\thepage}
  \setheadrule{0.35pt}}

\newpagestyle{bookhead}[\small\sffamily\color{bookslate!82!black}]{
  \sethead[\BookTitle][][\CTEXthechapter\quad\chaptertitle]
          {\CTEXthechapter\quad\chaptertitle}{}{\BookTitle}
  \setfoot[\thepage][][]
          {}{}{\thepage}
  \setheadrule{0.35pt}}

\renewpagestyle{plain}[\small\sffamily\color{black!58}]{
  \sethead{}{}{}
  \setfoot[][\thepage][]{}{\thepage}{}}

\definecolor{noteblue}{RGB}{238,244,247}
\definecolor{codegray}{RGB}{246,246,244}

% 章首插图不修改原图的语义内容，只在成书装配时用纸白覆盖层做四边渐隐。
% 图的尺寸由构建器原有的 adjustbox 决定；此宏不缩放、不裁切图片。
\tikzfading[name=chapterart-fade-west,
  left color=transparent!0,right color=transparent!100]
\tikzfading[name=chapterart-fade-east,
  left color=transparent!100,right color=transparent!0]
\tikzfading[name=chapterart-fade-south,
  bottom color=transparent!0,top color=transparent!100]
\tikzfading[name=chapterart-fade-north,
  bottom color=transparent!100,top color=transparent!0]
\newcommand{\chapterartfade}[1]{%
  \begin{tikzpicture}[baseline=(chapterart.base)]
    \node[inner sep=0pt,outer sep=0pt] (chapterart) {#1};
    \begin{scope}
      \clip (chapterart.south west) rectangle (chapterart.north east);
      \fill[white,path fading=chapterart-fade-west]
        (chapterart.south west) rectangle
        ($(chapterart.north west)!0.10!(chapterart.north east)$);
      \fill[white,path fading=chapterart-fade-east]
        ($(chapterart.south east)!0.10!(chapterart.south west)$) rectangle
        (chapterart.north east);
      \fill[white,path fading=chapterart-fade-south]
        (chapterart.south west) rectangle
        ($(chapterart.south east)!0.12!(chapterart.north east)$);
      \fill[white,path fading=chapterart-fade-north]
        ($(chapterart.north west)!0.12!(chapterart.south west)$) rectangle
        (chapterart.north east);
      % 覆住透明渐变在部分 PDF 渲染器中留下的单像素图片边界。
      \draw[white,line width=1.2pt]
        (chapterart.south west) rectangle (chapterart.north east);
    \end{scope}
  \end{tikzpicture}%
}

\newtcolorbox{noteblock}{enhanced,colback=noteblue,frame hidden,
  borderline west={2pt}{0pt}{bookteal!65!bookslate},
  left=7pt,right=7pt,top=5pt,bottom=5pt,breakable}
\newtcolorbox{codebox}{colback=codegray,colframe=black!25,
  boxrule=0.4pt,left=6pt,right=6pt,top=2pt,bottom=2pt,breakable}
\newtcolorbox{introbox}{enhanced,colback=bookamber!8,frame hidden,
  borderline west={2.2pt}{0pt}{bookamber!85!black},
  left=9pt,right=8pt,top=7pt,bottom=7pt,breakable}

\setcounter{tocdepth}{1}
\setcounter{secnumdepth}{2}

% 图号采用连字号式“图 X-M”（附录为“图 A-M”），与正文占位编号口径一致。
\renewcommand{\thefigure}{\thechapter-\arabic{figure}}

% 线框图 PDF（rsvg-convert 产出）为 1.7 版；抬升输出版本消除 xdvipdfmx 告警。
\AtBeginDocument{\special{pdf:minorversion 7}}
